\section{Uniform estimates for the matrix transfer density}
\label{ang:section}\label{sec:angular}

The previous section constructs a smooth matrix density on every sphere
of positive radius.  We prove that its $L^2$ norm stays bounded when the
spheres shrink, provided the two normalized principal matrices differ by
$O(r)$.  The main analytic input is a scalar coercivity argument from
\citet[Section~5]{ScalarCompanion2026}, whose proof we reproduce below.
We then compare the matrix equation with that scalar operator.  The
comparison is of lower differential order; it requires bounded matrix
drifts, but no smallness assumption on their difference.

\subsection{Moving traces and the matrix equation}

Use the reference metric and normalized inverse matrices from
\eqref{sph:normalization}.  The harmonic-frame convention
\eqref{bdy:harmonic-normalization}, with the multipliers $c_j$ in
\eqref{sph:frame-normalization}, gives equations whose second-derivative
coefficient matrices are positive definite:
\begin{equation}\label{ang:frame-equations}
 \mathcal L_1=a^{ij}\partial_i\partial_j+B_1^i\partial_i,
 \qquad
 \mathcal L_2=Q^{ij}\partial_i\partial_j+B_2^i\partial_i.
\end{equation}
The $B_j^i$ are complex $2\times2$ matrices acting on columns from the
left.  There is no zeroth-order term because each column of the defining
frame $F_j$ is harmonic.  All coefficients have uniform smooth bounds
on the fixed coordinate neighborhood for sufficiently small spatial
dilations.  Suppress the dilation and shrinking parameters for now, and
write
\[
 r=r_\delta,\qquad f_*=\log r,\qquad
 \beta=|df_*|_\gamma,\qquad f_i=\nabla_i f_*.
\]
The profile bounds supplied by Lemma~\ref{sph:profiles} have the form
\begin{equation}\label{ang:profile-bounds}
 \beta\ge\beta_0,\qquad
 \operatorname{Hess}_\gamma f_*\le-c_0\beta\gamma,\qquad
 |\nabla^j f_*|\le C_j\beta^j\quad(1\le j\le8).
\end{equation}
The lower threshold $\beta_0$ can be increased by increasing $p_*$.
Throughout this section we impose
\begin{equation}\label{ang:smallness}
 \sum_{|\alpha|\le8}|\partial_y^\alpha h|\le\varepsilon r,
 \qquad r\beta^2\le\varepsilon.
\end{equation}
The first is the bootstrap assumption of Section~\ref{sec:spheres};
the second is ensured by the choice of radius amplitude.

Let $S_\gamma Y$ be the tangent unit-sphere bundle with fiber probability
measure $d\sigma_y$.  Horizontal differentiation $\nabla$ uses parallel
transport for $\gamma$, including any base tensor indices, and acts
componentwise on the fixed harmonic-frame components.  The measure
$d\sigma_y$ is parallel.  Our Hilbert norms use
$dV_\gamma(y)d\sigma_y(\omega)$ and the standard Hermitian norm, or the
Frobenius norm for matrices.

There are two different transpose operations in the argument.  If
$\mathcal C$ is a sphere operator on columns, its \emph{bilinear row
transpose} $\mathcal C^t$ is defined by
\begin{equation}\label{ang:row-transpose}
 \int_{S_y}(\mathcal C^t\psi)V\,d\sigma_y
       =\int_{S_y}\psi(\mathcal C V)\,d\sigma_y.
\end{equation}
Here the identity is applied to each row of $\psi$; it involves no
complex conjugation.  In particular, a matrix multiplication $M$ on
columns becomes right multiplication $\psi\mapsto\psi M$, whereas a
matrix multiplying the output of the transfer map acts on $\psi$ from
the left.  A star will denote the Hermitian sphere adjoint of a scalar
operator, and a dagger its full Hilbert-space adjoint, including the
base integration.  Real scalar sphere operators have the same
bilinear transpose and Hermitian adjoint.

For a column $v$ solving $\mathcal L_1v=0$ near a closed moving ball, set
\[
 V(y,\omega)=v\bigl(\exp_y^\gamma(r(y)\omega)\bigr).
\]
Solving the ball Dirichlet problem expresses the boundary derivatives of
$v$ in terms of its trace and therefore defines sphere operators
$\mathcal C_i$ satisfying $\nabla_iV=\mathcal C_iV$.
For scalar functions write $\Delta_g=\operatorname{div}_g\nabla$.
Let $\mathcal C_{i,0}$ be the corresponding moving trace operators for
$c_1\Delta_{g_1}$, whose second-derivative coefficient matrix is
$a=\gamma^{-1}$.  This scalar comparison equation need not have zero
coordinate drift.

On each tangent unit ball, let $\mathsf B$ be the radial boundary
derivative of Euclidean harmonic extension and $\mathsf A_i$ its ambient
derivatives in a $\gamma$-orthonormal frame.  On spherical harmonics
$\mathcal H_k$ of degree $k$, $\mathsf B=k$ and
$\mathsf A_i:\mathcal H_k\to\mathcal H_{k-1}$.  The operators are parallel
with their tensor indices, and
\begin{equation}\label{ang:elementary-identities}
 [\mathsf B,\mathsf A_i^*]=\mathsf A_i^*,
 \qquad \sum_i\mathsf A_i^*\mathsf A_i^*=0.
\end{equation}
The second identity is the adjoint of the vanishing Laplacian of a
harmonic extension.  Define
\[
 \|w(y,\cdot)\|_s=\|(1+\mathsf B)^s w(y,\cdot)\|_{L^2(S_y)}.
\]
These norms are equivalent to the usual sphere Sobolev norms, uniformly
in $y$.  They also apply componentwise to tensors and matrices.

Write $\Psi^m$ for classical angular pseudodifferential operators of
order $m$, scalar or matrix valued as specified.  The notation
$c(y)\Psi^m$ means that division by $c(y)>0$ leaves uniformly bounded
symbol and operator seminorms in the finitely many orders being used.
An allowance $\beta^j$ for $j$ base derivatives means the corresponding
seminorms are bounded by $C_j\beta^j$.

\begin{lemma}[Trace expansions]\label{ang:trace}
Under \eqref{ang:profile-bounds}--\eqref{ang:smallness},
\begin{equation}\label{ang:C-expansion}
 \mathcal C_{i,0}^*
 =r^{-1}\mathsf A_i^*+f_i\mathsf B+E_i+F_i,
 \qquad E_i\in r\Psi^1,\quad F_i\in\Psi^0.
\end{equation}
The normalized remainders $r^{-1}E_i,F_i$ have the allowance $C_j\beta^j$
through six base derivatives.  Moreover, on rows,
\begin{equation}\label{ang:trace-difference}
 K_i:=\mathcal C_i^t-\mathcal C_{i,0}^*
       \in\Psi^0,\qquad \nabla K_i\in\beta\Psi^0.
\end{equation}
All constants depend on fixed smooth background bounds, not on the
lower threshold $\beta_0$ once $\beta_0\ge1$.
\end{lemma}

\begin{proof}
First regard the radius $s$ as an independent parameter.  Pull the first
scalar equation back by $z\mapsto\exp_y^\gamma(sz)$ and multiply by
$s^2$.  Its principal coefficients are $\delta^{ij}+O(s^2)$ and its
first-order coefficients are $O(s)$.  The same statement holds for the
system, with scalar principal coefficients times $I$.  Both Dirichlet
operators extend smoothly to $s=0$, where they equal the Euclidean
Laplacian.  On each fixed Sobolev scale their Dirichlet inverses are
uniformly bounded and smooth for sufficiently small $s$, by elliptic
Dirichlet estimates and invertibility at zero.

For these smooth strongly elliptic Dirichlet families on the unit ball,
the boundary radial derivative operator is classical of order one.  Its
principal symbol depends only on the principal coefficients and the
differentiating vector.  The boundary symbol construction, obtained by
choosing the decaying normal root and solving the successive transport
equations, is smooth in $(y,s)$; see, for example,
\citet[Section~3, Proposition~3.1]{JoshiLionheart2005}.
Their factorization is for the Hodge Laplacian; its construction applies
to the present Dirichlet systems because their
principal symbol is scalar times $I$, while the transport equations
allow matrix coefficients.  The true Dirichlet inverse corrects the
parametrix by a smoothing operator.  Differentiating this inverse on
arbitrarily high fixed Sobolev scales bounds every required smooth
kernel seminorm of the correction.  Thus these assertions hold in the
full symbol topology, including smoothing remainders and parameter
derivatives.

At $s=0$ the radial operator is $\mathsf B$, and the first $s$-derivative
of its principal symbol vanishes because the principal coefficients
have no linear term in $s$.  Taylor expansion in the symbol topology
therefore gives, for the scalar equation,
\begin{equation}\label{ang:radial-expansion}
 \mathsf B+s\mathcal B_1(y)+s^2\mathcal B_2(y,s),
 \qquad \mathcal B_1\in\Psi^0,\quad\mathcal B_2\in\Psi^1.
\end{equation}
The system and scalar radial operators have identical principal
symbols for every $s$ and coincide at zero.  Their difference is
consequently $s$ times a smooth family in $\Psi^0$.

Differentiating the center of the exponential map while transporting
$\omega$ parallel gives a Jacobi field with initial derivative zero.
In the unit-ball coordinates its velocity is
\[
 s^{-1}(e_i+O(s^2));
\]
putting $s=r(y)$ adds $f_i\omega$.  Tangential derivatives act directly
on boundary data and are identical for the scalar and matrix equations.
For the radial component use \eqref{ang:radial-expansion}.  The leading
terms are $r^{-1}\mathsf A_i+f_i\mathsf B$; the order-one errors have
size $O(r)$ and the order-zero errors are bounded, since
$r\beta\le\varepsilon/\beta$.  This proves \eqref{ang:C-expansion}
after sphere adjunction.  The radial difference $r\Psi^0$, multiplied
by the center velocity $O(r^{-1}+\beta)$, proves
\eqref{ang:trace-difference}.  Base derivatives cost at most one factor
$\beta$ each by the profile bounds.  Sphere transposition and
commutation with the smooth angular connection fields preserve these
orders and bounds.
\end{proof}

Write the second system in the reference connection as
$\mathcal L_2=Q^{ij}\nabla_i\nabla_j+b_2^i\nabla_i$; equivalently,
$b_2^k=B_2^k+Q^{ij}\Gamma_{ij}^k(\gamma)I$ in coordinates.
Set $D_i=\nabla_i+\mathcal C_i^t$.  The matrix density of
Proposition~\ref{sph:transfer} satisfies
\begin{equation}\label{ang:P-equation}
 P\psi=0,\qquad
 P=r\left[Q^{ij}D_iD_j+b_2^iD_i\right].
\end{equation}
Products are covariant products, and $b_2^i$ multiplies $D_i\psi$ from
the left.  In particular, $Q$ is not differentiated by a base adjoint.
Indeed, differentiation of $T_yv=\int\psi V\,d\sigma_y$ gives
$\nabla_i(T_yv)=\int(D_i\psi)V\,d\sigma_y$ by
\eqref{ang:row-transpose}.  Differentiate once more and apply the second
system equation, which holds for every local first-system harmonic
$v$ by Proposition~\ref{sph:transfer}.  The resulting expression pairs
to zero with all such traces.  These traces are dense in smooth boundary
data: solve the Dirichlet problems on slightly enlarged balls with
transported smooth boundary values and let their radii decrease to the
given radius.  Smooth parameter dependence gives convergence in every
fixed boundary smooth norm.  This proves \eqref{ang:P-equation}.

Write the scalar comparison equation on the second side as
$c_2\Delta_{g_2}=Q^{ij}\nabla_i\nabla_j+b_{2,0}^i\nabla_i$.
Define the scalar operator,
acting entrywise on matrices,
\begin{equation}\label{ang:P0}
 P_0=r\left[Q^{ij}D_{i,0}D_{j,0}+b_{2,0}^iD_{i,0}\right],
 \qquad D_{i,0}=\nabla_i+\mathcal C_{i,0}^*.
\end{equation}
The next subsections prove coercivity for $P_0$.  The proof uses
only its trace expansion and coefficient bounds; it assumes no scalar
transfer density or scalar Green-kernel identity.

\subsection{The model operator and its principal perturbations}

We separate a degree-raising model operator from the remaining
principal and lower-order terms of the moving trace equation.  We continue under the hypotheses of
Lemma~\ref{ang:trace}, including \eqref{ang:smallness}.

Define
\begin{equation}\label{ang:T-definition}
 D_i^+=\nabla_i+f_i\mathsf B,\qquad D_i^-=-\nabla_i+f_i\mathsf B,
 \qquad
 T=\sum_i \mathsf A_i^*D_i^+,\quad T^\flat=\sum_i \mathsf A_iD_i^-.
\end{equation}
The symbol $T$ here denotes a differential operator on sphere functions;
the transfer functional continues to be denoted $T_y$.
Expand \eqref{ang:P0} by
\eqref{ang:C-expansion}.  The identities
\eqref{ang:elementary-identities} give
\[
 D_i^+(r^{-1}\mathsf A_j^*)=r^{-1}\mathsf A_j^*D_i^+,
 \qquad \sum_i\mathsf A_i^*\mathsf A_i^*=0.
\]
Thus, in ordinary base coordinates and an orthonormal sphere
trivialization,
\begin{equation}\label{ang:P-decomposition}
 P_0=2T+\mathcal U+\mathcal J,
\end{equation}
where $\mathcal J$ is a uniformly bounded family in $\Psi^1$, with no base
derivatives, and
\begin{equation}\label{ang:U-bounds}
 \mathcal U=\sum_{|\alpha|\le2}\mathcal U_\alpha(y)\partial_y^\alpha,
 \qquad \mathcal U_\alpha\in\varepsilon\beta^{-|\alpha|}
                               \Psi^{2-|\alpha|}.
\end{equation}
Up to five derivatives of these coefficients have the additional
allowance $C_j\beta^j$.  The second-base-derivative term is precisely
\begin{equation}\label{ang:real-principal-part}
 rQ^{ij}\partial_{y_i}\partial_{y_j};
\end{equation}
its coefficients are real scalars independent of the sphere variable.

We detail these bounds.  The pure angular second derivative contracted
with $a$ vanishes.  The remaining contraction is
$(h^{ij}/r)\mathsf A_i^*\mathsf A_j^*$, and belongs to the zero-base-derivative term of
\eqref{ang:U-bounds}.  The mixed correction is
$2h^{ij}\mathsf A_i^*D_j^+$ and has the stated small bounds.  The $D^+D^+$ term
has coefficient $r$; its base orders two, one and zero have coefficients
$O(r)$, $O(r\beta)$ and $O(r\beta^2)$, respectively.  The condition
$r\beta^2\le\varepsilon$ gives \eqref{ang:U-bounds} for all three.
Terms containing $E_i$ obey the same bounds.  Terms containing $F_i$
also do so, except for their cross terms with $r^{-1}\mathsf A_j^*$, which are
bounded angular order-one terms and belong to $\mathcal J$.  The drift term with
$\mathsf A_i^*$ belongs to $\mathcal J$ as well.  Differentiating any of these expressions
gives the stated coefficient bounds by
\eqref{ang:profile-bounds} and \eqref{ang:smallness}.
A horizontal connection written in this trivialization adds a bounded
angular first derivative, which preserves the same estimates.  The bounds
for $\mathcal J$ use the fixed background and remain uniform as $\beta_0$ is
increased.

The decomposition identifies the model operator $2T$ and the
second-order perturbation $\mathcal U$.  To carry the model estimate over to $P_0$
we shall need more than an ordinary first-derivative estimate: the error forms
also contain mixed base and angular derivatives.  For a sphere function
$\varphi$, put
\begin{equation}\label{ang:N-definition}
 \mathcal N(\varphi)^2
 =\int_Y\left(
  \|\nabla\varphi\|_0^2+\beta^2\|\varphi\|_1^2
  +\beta^{-1}\|\nabla\varphi\|_{1/2}^2
  +\beta\|\varphi\|_{3/2}^2\right)dV_\gamma.
\end{equation}

The last two terms of $\mathcal N$ place half an angular derivative
on $\nabla\varphi$ and three halves on $\varphi$.  They will bound the
remaining error forms left by the comparison with $P_0$.

\subsection{A comparison of raising and lowering operators}

The next estimate supplies these two derivative strengths for the model
operator.  Spherical harmonics are trace-free symmetric tensors, and
comparing their raising and lowering operators is familiar from energy
identities in tensor tomography; see, for example, \cite{PSU15,GPSU16}.
Here the negative Hessian of the spatial weight supplies the positive
term that absorbs the bounded curvature.  The particular weighted estimate below is adapted from
\citet[Section~5, weighted raising estimate]{ScalarCompanion2026}.
We include the degree calculation because it is the source of the
additional half angular derivative used in the perturbation argument.

\begin{lemma}[Weighted raising estimate]\label{ang:raising}
Let $n\ge3$, and suppose \eqref{ang:profile-bounds} holds on a fixed
base region with uniformly bounded reference geometry.  There are
$\beta_0$ and $c>0$, depending only on that geometry and the constants in
\eqref{ang:profile-bounds}, such that every smooth, compactly
base-supported function with zero fiber mean satisfies
\begin{equation}\label{ang:raising-estimate}
 \|T\varphi\|^2-\|T^\flat\varphi\|^2
       \ge c\,\mathcal N(\varphi)^2.
\end{equation}
\end{lemma}

\begin{proof}
Both operators change angular degree by exactly one, so it suffices to
prove the estimate at one input degree $k\ge1$.  Identify that component
with a harmonic homogeneous polynomial $v(z)$ of degree $k$.  Use the
Fischer inner product, in which the monomials $z^\alpha$ are orthogonal
with squared norm $\alpha!$.  Polynomial differentiation $a_i=\partial_{z_i}$
is adjoint to multiplication $M_i=z_i$.  On harmonic polynomials of degree
$k$ the Fischer squared norm is the normalized squared sphere norm
multiplied by
\[
 N_k=n(n+2)\cdots(n+2k-2),\qquad N_0=1.
\]
For completeness, integration against a standard real Gaussian identifies
the Fischer norm on trace-free homogeneous polynomials with their
Gaussian $L^2$ norm;
integrating radially then gives the factor $N_k$.  These classical
identities are also recorded in \cite[Equation~(2.1) and
Theorem~2.1]{Render08}.

Set $d=k+n/2-1$ and use $D_i^\pm=\pm\nabla_i+k f_i$ at this fixed degree.
For harmonic polynomials $q_i$ of degree $k$, the Laplacian of
$\sum_iM_iq_i$ is $2\sum_i a_iq_i$.  Since
\[
 \Delta_z(|z|^2s)=4d\,s\qquad(s\in\mathcal H_{k-1}),
\]
its harmonic projection is
$\sum_iM_iq_i-|z|^2\sum_i a_iq_i/(2d)$.  Orthogonality of this projection
and its complement shows that its squared norm is
\[
 \sum_i\|q_i\|^2+\sum_{i,j}(a_jq_i,a_iq_j)
             -d^{-1}\Bigl\|\sum_i a_iq_i\Bigr\|^2.
\]
Changing from Fischer adjoints to sphere adjoints gives
\[
 \mathsf A_i^*|_{\mathcal H_k}
   =\frac{N_{k+1}}{N_k}\Pi_{k+1}M_i,
\]
where $\Pi_{k+1}$ is harmonic projection in homogeneous degree $k+1$.
Consequently $N_k$ times the left side of
\eqref{ang:raising-estimate} is
\begin{equation}\label{ang:Fock-difference}
 \begin{split}
 (2d+2)\left[\|D^+v\|^2
   +\sum_{i,j}(a_jD_i^+v,a_iD_j^+v)
   -\frac1d\Bigl\|\sum_i a_iD_i^+v\Bigr\|^2\right]
 -2d\Bigl\|\sum_i a_iD_i^-v\Bigr\|^2.
 \end{split}
\end{equation}
All norms in this calculation include integration over the base and use
the polynomial inner product in the fibers.  The factors are
$N_{k+1}/N_k=2d+2$ and $N_k/N_{k-1}=2d$.

Write $C_{ij}=2k(\operatorname{Hess}_\gamma f_*)_{ij}$ and
\[
 C[av]=\int_Y\sum_{i,j}C_{ij}\langle a_iv,a_jv\rangle\,dV_\gamma.
\]
Covariant integrations by parts give
\begin{align}
 \sum_{i,j}(a_jD_i^+v,a_iD_j^+v)
   &=\Bigl\|\sum_i a_iD_i^-v\Bigr\|^2-C[av]
                                    +O(k^2)\|v\|^2,\label{ang:ibp-one}\\
 \Bigl\|\sum_i a_iD_i^+v\Bigr\|^2
   &\le k\|D^-v\|^2-C[av]+O(k^2)\|v\|^2,\label{ang:ibp-two}\\
 \|D^-v\|^2
   &=\|D^+v\|^2+\int_Y\tr C\,|v|^2\,dV_\gamma.
                                                    \label{ang:ibp-three}
\end{align}
Here and below a form denoted $O(k^2)\|v\|^2$ has absolute value at most
a fixed constant times that quantity.  To verify the first two formulas,
move a $D_j^+$ across the pairing and commute
$D_j^-D_i^+$ to $D_i^+D_j^--C_{ij}$.  The reordered term in
\eqref{ang:ibp-two} is at most $k\|D^-v\|^2$, because
$\sum_i|a_iw|^2=k|w|^2$ for $w\in\mathcal H_k$.
The extra commutator is curvature: it acts as a sum of rotations of norm
$O(k)$ at degree $k$, and the two contractions supply one more factor $k$.
This proves the claimed error bounds.  Expanding the two squares proves
\eqref{ang:ibp-three}.  For complex functions all the displayed forms
are understood through their real parts.

Substituting \eqref{ang:ibp-one} into
\eqref{ang:Fock-difference} rewrites that expression as
\begin{equation}\label{ang:after-first-ibp}
 \begin{split}
 &(2+2/d)\left[d\|D^+v\|^2
       -\Bigl\|\sum_i a_iD_i^+v\Bigr\|^2-dC[av]\right]\\
 &\qquad+2\Bigl\|\sum_i a_iD_i^-v\Bigr\|^2
                      +O(k^3)\|v\|^2.
 \end{split}
\end{equation}
We must control the remaining negative divergence square while retaining
both an ordinary derivative estimate and an additional half-order
angular derivative estimate.  The latter requires a coefficient of
size $k/\beta$ in front of $|D^+v|^2$, including when $k$ is much larger
than $\beta$.  Choose a small constant $\vartheta>0$.
On the fraction $1-\vartheta$ of its negative divergence square use
\eqref{ang:ibp-two}--\eqref{ang:ibp-three}.  On the remaining
fraction use the pointwise convex combination
\begin{equation}\label{ang:convex-bound}
 \begin{split}
 \Bigl|\sum_i a_iD_i^+v\Bigr|^2
 &\le(1-\zeta)(k+n-1)|D^+v|^2\\
 &\quad+\zeta\left(2\Bigl|\sum_i a_iD_i^-v\Bigr|^2
                         +8k^3\beta^2|v|^2\right),
 \qquad \zeta=c_1/\beta.
 \end{split}
\end{equation}
The fraction $1-\vartheta$ in the integrated estimates is constant;
the $y$-dependent $\zeta$ is used only in this pointwise inequality.
The first bound follows from the adjoint multiplication operators, since
the contraction $(q_i)\mapsto\sum_i a_iq_i$ has squared norm at most
$k+n-1$ on polynomials of degree $k$.  The second follows from
$D_i^++D_i^-=2k f_i$ and
$|\sum_i f_i a_i v|^2\le k\beta^2|v|^2$.

Combining these bounds in \eqref{ang:after-first-ibp}, the coefficient
retained for $|D^+v|^2$ is
\begin{equation}\label{ang:positive-degree-factor}
 (2+2/d)\bigl[d-k-\vartheta(n-1)
                         +\vartheta\zeta(k+n-1)\bigr].
\end{equation}
Since $d-k=(n-2)/2$, a sufficiently small fixed $\vartheta$ makes this
at least a fixed positive constant plus a positive multiple of
$c_1k/\beta$.  The coefficient of
$|\sum_i a_iD_i^-v|^2$ is nonnegative when $\beta_0$ is sufficiently
large.  The Hessian contribution is exactly
\begin{equation}\label{ang:hessian-contribution}
 -(2+2/d)\left[(d-1+\vartheta)C[av]
       +(1-\vartheta)k\int_Y\tr C\,|v|^2\,dV_\gamma\right].
\end{equation}
For $k\ge1$ and $n\ge3$, both coefficients inside this expression are
positive.  The Hessian hypothesis bounds it below by
$c k^3\int\beta|v|^2$.  This absorbs the curvature error $O(k^3)\|v\|^2$
by increasing $\beta_0$, and absorbs the last term of
\eqref{ang:convex-bound} by choosing $c_1$ sufficiently small.
We have therefore retained positive multiples of
\[
 \|D^+v\|^2,\qquad
 \int_Y\frac{k}{\beta}|D^+v|^2\,dV_\gamma,
 \qquad k^3\int_Y\beta|v|^2\,dV_\gamma.
\]
The first controls both $\|\nabla v\|^2$ and
$k^2\int\beta^2|v|^2$, since its cross term is
$-k\int\Delta_\gamma f_*\,|v|^2\ge0$.  For the weighted derivative use the
pointwise inequality
\[
 \frac{k}{\beta}|\nabla v|^2
 \le2\frac{k}{\beta}|D^+v|^2+2k^3\beta|v|^2.
\]
No differentiation of $\beta^{-1}$ is required.  Dividing by $N_k$, and
then summing the orthogonal degree components, proves
\eqref{ang:raising-estimate}; the weights $(1+k)^s$ are comparable
to $k^s$ because the zero degree is excluded.
\end{proof}

\subsection{Keeping the principal perturbations}

The raising estimate controls the model operator.  The scalar operator
\eqref{ang:P0} also contains second base derivatives and a
small second-order angular perturbation.  We retain these terms in a
comparison of squared norms.

\begin{proposition}[Scalar coercivity]
\label{ang:scalar-coercivity}
Fix uniform smooth bounds for the reference geometry and scalar
comparison equations, and the constants in \eqref{ang:profile-bounds}.
Suppose $P_0$ is the scalar operator
\eqref{ang:P0}, with the scalar trace operators of
Lemma~\ref{ang:trace}.  There exist $\beta_0$ sufficiently large,
$\varepsilon_0>0$ and $C<\infty$, depending only on these fixed bounds, such that
\eqref{ang:smallness} with $\varepsilon\le\varepsilon_0$ implies
\begin{equation}\label{ang:coercive-estimate}
 \mathcal N(\varphi)\le C\|P_0\varphi\|
\end{equation}
for every smooth, compactly base-supported function $\varphi$ with
zero mean on each sphere.  These constants are independent of the
shrinking parameter.  Only the eight indicated derivatives of the small
tensor $h$ are required.
\end{proposition}

\begin{proof}
Let a dagger denote the full adjoint for $dV_\gamma d\sigma$.
Covariant integration and the degree shift give
\begin{equation}\label{ang:Z-definition}
 Z:=T^\flat-T^\dagger=\sum_i f_i\mathsf A_i.
\end{equation}
A direct expansion yields
\begin{equation}\label{ang:norm-comparison}
 \begin{split}
 &\|(2T+\mathcal U)\varphi\|^2-
                \|(2T^\flat+\mathcal U^\dagger)\varphi\|^2\\
 &\qquad=4\bigl(\|T\varphi\|^2-\|T^\flat\varphi\|^2\bigr)
                       +\langle\mathcal E\varphi,\varphi\rangle,
 \end{split}
\end{equation}
where
\begin{equation}\label{ang:error-operator}
 \mathcal E=2[T^\dagger,\mathcal U]+2[\mathcal U^\dagger,T]+[\mathcal U^\dagger,\mathcal U]
                       -2\bigl(\mathcal U Z+(\mathcal U Z)^\dagger\bigr).
\end{equation}
We will prove
\begin{equation}\label{ang:error-form-bound}
 |\langle\mathcal E\varphi,\varphi\rangle|
                         \le C\varepsilon\mathcal N(\varphi)^2.
\end{equation}

Here are the complete order and weight counts needed for this claim.
An operator has count $(m,s)$ if its coefficient at $\partial_y^\alpha$
is angular order $m-|\alpha|$, of size $O(\beta^{s-|\alpha|})$,
with the derivative allowances already specified.
Adjoints preserve the count, and products add the counts.  When a base
derivative hits a coefficient it consumes one differential order and
adds a factor $\beta$, so it preserves the weight index.  A commutator
of scalar angular operators loses one angular order: the products of
their principal symbols cancel.  Hence a scalar commutator loses one in
total order while preserving the sum of the weight indices.
The counts of $T,T^\dagger,\mathcal U,Z$ are, respectively,
\[
 (2,1),\quad(2,1),\quad(2,0),\quad(1,1),
\]
and every occurrence of $\mathcal U$ carries the factor $\varepsilon$.  Thus
\eqref{ang:error-operator} has count at most $(3,1)$ and a factor
$O(\varepsilon)$; the quadratic $\mathcal U$ term is smaller.

There is a further cancellation that the total count alone does not
express: no third base derivative survives.  To check it explicitly,
use an orthonormal sphere trivialization and write the base measure as
$\mu(y)\,dy$.  The sphere probability measure is then fixed.  Write
\[
 \mathcal U=c^{ij}\partial_i\partial_j+\mathcal B^i\partial_i+\mathcal C,
 \qquad c^{ij}=rQ^{ij}=c^{ji}\in\mathbb R.
\]
The coefficients $\mathcal B^i$ and $\mathcal C$ are angular
operators.  Replacing a horizontal derivative by a coordinate derivative
plus its angular connection term changes only base orders at most one,
so these are all the second-base coefficients.  If a star denotes the
sphere adjoint, direct integration by parts gives
\begin{align*}
 \mathcal U^\dagger-\mathcal U
 ={}&\left(2\mu^{-1}\partial_j(\mu c^{ij})
                -\mathcal B^i-(\mathcal B^i)^*\right)\partial_i\\
 &+\mu^{-1}\partial_i\partial_j(\mu c^{ij})
       -\mu^{-1}\partial_i\bigl(\mu(\mathcal B^i)^*\bigr)
       +\mathcal C^*-\mathcal C.
\end{align*}
In particular it has base order at most one.
In $[T^\dagger,\mathcal U]$ and
its partner, the possible third-base-derivative coefficient is a commutator of an angular coefficient
with $rQ^{ij}$, which is zero by
\eqref{ang:real-principal-part}.  For $[\mathcal U^\dagger,\mathcal U]$, first write
it as $[\mathcal U^\dagger-\mathcal U,\mathcal U]$.  The operator $\mathcal U^\dagger-\mathcal U$ has no second base
derivative, since the coefficient in
\eqref{ang:real-principal-part} is real.  Taking adjoints against
the smooth base density only produces lower base orders.  Its remaining
possible third-base-derivative coefficients commute for the same reason.
Finally, $\mathcal U Z$ already has at most two base derivatives.
It follows that $\mathcal E$ is a sum of terms of precisely the following
three permitted types:
\begin{equation}\label{ang:error-types}
 \varepsilon\beta^{-1}\Psi^1\partial_y^2,
 \qquad \varepsilon\Psi^2\partial_y,
 \qquad \varepsilon\beta\Psi^3.
\end{equation}
This argument also covers lower orders, since $\beta\ge1$.

For the first type in \eqref{ang:error-types}, integrate once in the
base.  The principal form is bounded by
$C\varepsilon\int\beta^{-1}\|\partial_y\varphi\|_{1/2}^2$.
A derivative of its coefficient costs one factor $\beta$ and leaves an
$\varepsilon\Psi^1\partial_y$ term.  This includes differentiation of the
weight itself: $|d\beta|\le C\beta^2$ implies
$|d(\beta^{-1})|\le C$.  The second type has the form bound
\[
 C\varepsilon\int_Y
     \|\partial_y\varphi\|_{1/2}\|\varphi\|_{3/2}\,dV_\gamma
 \le C\varepsilon\int_Y\left(
     \beta^{-1}\|\partial_y\varphi\|_{1/2}^2
                 +\beta\|\varphi\|_{3/2}^2\right)dV_\gamma.
\]
The third type is bounded by
$C\varepsilon\int\beta\|\varphi\|_{3/2}^2$.
The lower-order term from differentiating a coefficient satisfies the
same bound.  Replacing coordinate derivatives by horizontal derivatives
adds a bounded angular first derivative, again controlled by these norms.
A fixed finite partition of the base region has the same harmless
lower-order errors.  This proves \eqref{ang:error-form-bound}.

The constants in the preceding estimates are uniform for all
$\beta\ge\beta_0$ once a fixed lower threshold is met.  Choose
$\varepsilon_0$ small relative to the constant in
Lemma~\ref{ang:raising}.  Dropping the nonnegative second squared
norm in \eqref{ang:norm-comparison} gives
$\mathcal N(\varphi)\le C\|(2T+\mathcal U)\varphi\|$.
Finally,
\[
 \|\mathcal J\varphi\|\le C\|\varphi\|_{L^2(Y;H^1(S_y))}
                         \le C\beta_0^{-1}\mathcal N(\varphi),
\]
so a sufficiently large $\beta_0$ absorbs $\mathcal J$ and proves
\eqref{ang:coercive-estimate}.

All these operations use angular symbol estimates pointwise in the base,
and ordinary base differential operators of degree at most two.  Formal
adjoints, the commutator expansion and the one integration by parts use
at most the five base coefficient derivatives specified in
\eqref{ang:U-bounds}.  Forming these coefficients uses at most one base
derivative of a trace operator; six trace derivatives therefore suffice.
Since its weight dependence is through $r$ and $f_i$, profile derivatives
through order seven suffice, within the eight allowed in
\eqref{ang:profile-bounds}.  Differentiating a factor such as $h/r$ five
times uses only five derivatives of $h$, with the remaining factors
bounded by the permitted powers of $\beta$.  Higher angular symbol
seminorms depend on the fixed smooth first-metric family: the tensor $h$
is evaluated at the center, so angular differentiation introduces no
further center derivatives of it.  The scalar drifts have fixed smooth
background bounds.  Thus the eight smallness derivatives in
\eqref{ang:smallness} suffice in every fixed dimension; no smallness of
all derivatives of $h$ is used.
\end{proof}

\subsection{The matrix comparison and the density bound}

The scalar calculation is now complete.  We use it entrywise before
introducing any matrix coefficients into a squared-norm comparison.
This order matters: matrix principal symbols need not commute, whereas
the preceding commutator argument is scalar.

\begin{proposition}[System coercivity]\label{ang:system-coercivity}
Fix uniform smooth bounds for the harmonic-frame systems and the
reference geometry, and the constants in \eqref{ang:profile-bounds}.
There are $\beta_0$ sufficiently large, $\varepsilon_0>0$, and $C$,
depending only on these bounds, such that
\eqref{ang:smallness} with $\varepsilon\le\varepsilon_0$ implies
\[
 \mathcal N(\varphi)\le C\|P\varphi\|_{L^2}
\]
for every smooth compactly base-supported matrix function $\varphi$
with zero fiber mean.  The constants are independent of the shrinking
parameter and of sufficiently small spatial dilations.
\end{proposition}

\begin{proof}
Put $q^i=b_2^i-b_{2,0}^iI$ and use
$D_i=D_{i,0}+K_i$ from \eqref{ang:trace-difference}.
Let $L_M$ denote left multiplication by a matrix $M$.  Expansion, without
interchanging any factors, gives the identity
\begin{equation}\label{ang:comparison-identity}
 \begin{split}
 P-P_0=r\bigl[&Q^{ij}(D_{i,0}K_j+K_iD_{j,0}+K_iK_j)\\
              &+L_{q^i}D_{i,0}+L_{b_2^i}K_i\bigr].
 \end{split}
\end{equation}
Thus $D_{i,0}K_j$ denotes composition and includes differentiation of
$K_j$.  There is no derivative of $q^i$ or $b_2^i$ in this formula.
Left multiplication commutes with the row action of a sphere operator
when its matrix depends only on the center, but no such commutation is
needed in \eqref{ang:comparison-identity}.

The trace expansion gives
$\mathcal C_{i,0}^*\in(r^{-1}+\beta+1)\Psi^1$.
Together with $K_i\in\Psi^0$, $\nabla K_i\in\beta\Psi^0$, and bounded
drifts, \eqref{ang:comparison-identity} consequently yields
\begin{equation}\label{ang:comparison-bound}
 \|(P-P_0)\varphi\|_{L^2}
 \le C\left(\|\varphi\|_{L^2(Y;H^1(S_y))}
                         +\|r\nabla\varphi\|_{L^2}\right).
\end{equation}
To see the weights explicitly, every term with a derivative of
$\varphi$ in the base is $r\Psi^0\nabla\varphi$.
Products of $K_i$ with the leading $r^{-1}\mathsf A_j^*$ leave a bounded
$\Psi^1$ term after multiplication by $r$; products with
$f_j\mathsf B$ have coefficient $O(r\beta)$.
Differentiating $K_j$ gives $r\beta\Psi^0$.  The remaining products
have smaller orders and bounded coefficients.  Since $r\beta\le
\varepsilon/\beta$, all constants in \eqref{ang:comparison-bound}
are uniform as the lower threshold $\beta_0$ increases.

Apply Proposition~\ref{ang:scalar-coercivity} to the four entries of
$\varphi$.  By the definition of $\mathcal N$,
\[
 \|\varphi\|_{L^2(Y;H^1)}\le\beta_0^{-1}\mathcal N(\varphi),
 \qquad
 \|r\nabla\varphi\|_{L^2}
      \le\varepsilon\beta_0^{-2}\mathcal N(\varphi).
\]
It follows that
\[
 \mathcal N(\varphi)
 \le C_0\|P\varphi\|_{L^2}
   +C_0C(\beta_0^{-1}+\varepsilon\beta_0^{-2})
                                        \mathcal N(\varphi).
\]
Increase the fixed lower threshold so that the last coefficient is at
most $1/2$.  This proves the assertion.  In the parameter choices for
the construction, this increase is made when $p_*$ is chosen, before
$\varepsilon$ and the radius amplitude are fixed.
\end{proof}

We record two bounds needed to use coercivity on the actual density.
First, $P_0 1$ is uniformly bounded: the raising part in
\eqref{ang:P-decomposition} annihilates constants, the zeroth base
coefficient of $\mathcal U$ is bounded in $\Psi^2$, and $\mathcal J$ is bounded in
$\Psi^1$.  Equation~\eqref{ang:comparison-bound} then shows that
$PI$ is uniformly bounded in $L^2$.
Second, for any fixed smooth base cutoff $\chi$,
\begin{equation}\label{ang:cutoff-bound}
 \|[P,\chi]w\|_{L^2}
 \le C_\chi\left(
   \|w\|_{L^2(\mathcal O_\chi;H^1(S_y))}
             +\|r\nabla w\|_{L^2(S_\gamma\mathcal O_\chi)}\right),
\end{equation}
where $\mathcal O_\chi$ is any fixed small neighborhood of the supports
of its derivatives.  Indeed,
$[2T,\chi]=2\sum_i\mathsf A_i^*(\nabla_i\chi)$,
$\mathcal U$ has second base coefficient $rQ^{ij}$ and first base coefficients
bounded in $\Psi^1$, and $\mathcal J$ commutes with $\chi$.
The additional first base coefficients of $P-P_0$ are $r\Psi^0$,
by \eqref{ang:comparison-identity}.  These observations give exactly
\eqref{ang:cutoff-bound}.

\begin{proposition}[Uniform transfer density]\label{ang:density}
Choose the fixed geometry as in Section~\ref{sec:spheres}, then choose
$p_*$ sufficiently large, $\varepsilon>0$ sufficiently small, and the
radius amplitude $R_*>0$ sufficiently small that
$r_\delta\beta_\delta^2\le\varepsilon$ for $0<\delta\le1$.
For all sufficiently small spatial dilations there is a constant $C$,
independent of the dilation and $\delta$, such that
\begin{equation}\label{ang:density-bound}
 \|\chi\psi_\delta\|_{L^2(S_\gamma Y)}\le C
\end{equation}
whenever $\sum_{|\alpha|\le8}|\partial^\alpha h|
\le\varepsilon r_\delta$ on $Y$.
Here $\chi$ is the fixed cutoff from \eqref{sph:cutoff} and
$\psi_\delta$ is the left-acting matrix density of
Proposition~\ref{sph:transfer}.
\end{proposition}

\begin{proof}
The reproduction identity $\int\psi_\delta\,d\sigma_y=I$ makes
$\varphi=\chi(\psi_\delta-I)$ mean zero on each sphere.  It is smooth
and compactly supported in the base.  Equation~\eqref{ang:P-equation}
gives
\[
 P\varphi=-\chi PI+[P,\chi](\psi_\delta-I).
\]
The first term has the uniform bound just established.  For the second,
\eqref{ang:cutoff-bound} requires only one angular derivative and
$r\nabla\psi_\delta$ on the cutoff derivative region.
The uniform bounds of Lemma~\ref{sph:cutoff-bounds} supply these:
in the lower region the density is the normalized ball evaluation
kernel, and in the other region the radii have a fixed positive lower
bound.  System coercivity bounds $\mathcal N(\varphi)$, hence
$\|\varphi\|_{L^2}$, uniformly.  Adding the fixed matrix function
$\chi I$ proves \eqref{ang:density-bound}.
\end{proof}

The estimate controls the density of the actual transfer functional.
In the next section its first and second Taylor moments will turn this
bound into quantitative information about the difference of the two
harmonic-frame equations.
